\subsection{FRM EXAM 2009---QUESTION 1-11: rare event vs risk management failure}

\textbf{Enunciado}

\emph{Bank United’s CEO decided to make a large investment in a levered portfolio of CDOs. The CRO estimated that the portfolio had a 1\% chance of losing \$1 billion or more over one year. At the end of the first year the portfolio lost \$2 billion and the bank was closed. Which statement is correct?}
\begin{enumerate}[label=\alph*.]
\item The outcome demonstrates a risk management failure because the bank did not eliminate the possibility of financial distress.
\item The outcome demonstrates a risk management failure because an extremely unlikely outcome occurred, so the probability was poorly estimated.
\item The outcome demonstrates a risk management failure because the CRO failed to go to regulators to stop the shutdown.
\item \hl{Based on the information provided, one cannot determine whether it was a risk management failure.}
\end{enumerate}

\subsubsection*{Solución}


\textbf{Evento raro no implica error automáticamente. }Una probabilidad del 1\% significa que el evento puede ocurrir. Observar una sola pérdida extrema (\$2B) no prueba que la probabilidad estuviera mal estimada: se requeriría evidencia adicional (supuestos del modelo, backtesting, más observaciones, etc.).

\textbf{Riesgo residual y apetito de riesgo. } La gestión de riesgos normalmente no ``elimina'' toda posibilidad de insolvencia; busca cuantificar, limitar y capitalizar el riesgo según el apetito del banco. Sin más información, no se puede concluir fallo.


Con la información dada no se puede determinar, así que la opción (d) es la adecuada.












